\section{Chern--Simons 作用量与 $k$ 的整性}
\label{sec:cs}

\subsection{为什么不能直接用 Yang--Mills 作用量}

设 $M$ 是三维流形，$G$ 是紧单 Lie 群，$A=A_i\,\dd x^i$ 是 $M$ 上 $G$-联络（Lie 代数取值 1-形式）。规范变换的无穷小形式是 \src{353}
\begin{equation}
A_i\;\longrightarrow\; A_i-D_i\epsilon ,\qquad
F_{ij}=[D_i,D_j]=\partial_iA_j-\partial_jA_i+[A_i,A_j].
\label{eq:gauge}
\end{equation}

标准的 Yang--Mills 作用量是
\begin{equation}
L_0=\int_M\sqrt{g}\,g^{ik}g^{jl}\,\Tr(F_{ij}F_{kl}) .
\label{eq:ym}
\end{equation}
它显含度规 $g$，通过 $\sqrt g$ 与 $g^{ik}g^{jl}$ 两处 \src{353}。所以量子化后一般得不到不依赖度规的量。注意两个因子的性质不同：$\sqrt g\,\dd^3x$ 是体积元，换度规只改标度；但 $g^{ik}g^{jl}$ 把两个 $F$ 的指标缩并，换度规会改变不同分量之间的混合方式。前者可以靠“不积分体积”绕开，后者绕不开。

\subsection{Chern--Simons 作用量}

Witten 用的是 \src{354}（(1.3)）：
\begin{equation}
\begin{aligned}
L&=\frac{k}{4\pi}\int_M \Tr\Bigl(A\wedge \dd A+\frac23 A\wedge A\wedge A\Bigr)\\
 &=\frac{k}{8\pi}\int \dd^3x\,\epsilon^{ijk}\Tr\Bigl(A_i\partial_jA_k+\frac23 A_i[A_j,A_k]\Bigr).
\end{aligned}
\label{eq:cs}
\end{equation}

它的三个性质决定了一切：
\begin{enumerate}[leftmargin=2em]
\item \textbf{只依赖定向}：被积式是 3-形式，$M$ 的定向决定 $\epsilon^{ijk}$ 的符号；整个式子没有 $g$。
\item \textbf{变分给出曲率}：$\delta\int\Tr(A\dd A+\tfrac23A^3)=2\int\Tr(\delta A\wedge F)$，所以 \ours
\begin{equation}
\delta L=\frac{k}{2\pi}\int_M \Tr(\delta A\wedge F).
\label{eq:var}
\end{equation}
驻点条件是 $F=0$——平坦联络。这正是第 \ref{sec:canonical} 节相空间是“平坦联络模空间”的来源。
\item \textbf{大规范变换下不是不变，而是平移 $2\pi$ 的整数倍}：这是 $k$ 必须整数的根源，见下。
\end{enumerate}

\subsection{$k$ 的整性：大规范变换}

$\pi_3(G)\cong\ZZ$，所以规范变换 $g:M\to G$ 有绕数 $m$。在 $g$ 下作用量平移 \src{354}
\begin{equation}
L\;\longrightarrow\; L+\text{const}\cdot m .
\end{equation}
原文只写“常数乘以 $m$”，没有写出常数。我们把它算出来 \ours：对 $SU(N)$ 取 $N$ 维表示的迹，
\begin{equation}
\Delta L=\frac{k}{4\pi}\cdot\frac13\int \Tr(g^{-1}\dd g)^3
=\frac{k}{4\pi}\cdot\frac13\cdot 24\pi^2 m=2\pi k\,m ,
\label{eq:bigauge}
\end{equation}
其中用到 $\frac1{24\pi^2}\int\Tr(g^{-1}\dd g)^3\in\ZZ$。

量子理论要求 $\exp(\ii L)$ 单值，即 $\Delta L\in 2\pi\ZZ$，于是
\begin{equation}
k\in\ZZ .
\end{equation}
原文的说法是：这个条件“必要且充分” \src{355}。必要是上面这条；充分指的是：只要 $k$ 整，大规范变换就不改变 $\exp(\ii L)$。这是本文对“为什么耦合常数会量子化”的第一次回答——与二维仿射李代数中心荷的整性同源。

\begin{keypoint}
\textbf{约定陷阱}：$k$ 的整性与 $\Tr$ 的归一化绑定。若改换表示或改换 Killing 型的归一化（例如把长根平方取为 2），$k$ 的整性条件会跟着变。河野书用的是 $\mathrm{CS}=\frac1{8\pi^2}\int\Tr(A\dd A+\frac23A^3)$、配分函数 $\int e^{2\pi\ii k\,\mathrm{CS}}$ \src{河野 p.74}，与本文的 $L$ 相差一个常数因子；比较两书的 $k$ 时必须先对齐归一化。
\end{keypoint}

\subsection{Wilson 线与配分函数}

Wilson 线是沿闭曲线 $C$ 的路径序指数 \src{355}（(1.5)）：
\begin{equation}
W_R(C)=\Tr_R\,P\exp\int_C A_i\,\dd x^i .
\label{eq:wilson}
\end{equation}
它不需要度规：$A_i\dd x^i$ 是 1-形式沿曲线的积分，$\Tr_R$ 是表示空间的迹。

链环 $C_1,\dots,C_n$（带表示 $R_1,\dots,R_n$）的配分函数定义为 \src{355}（(1.6)）
\begin{equation}
Z(M;C_1,\dots,C_n)=\int \mathcal{D}A\;\; e^{\ii L}\prod_{i=1}^n W_{R_i}(C_i).
\label{eq:partition}
\end{equation}

本文的核心断言是：取 $M=S^3$、$G=SU(N)$、所有 $R_i$ 都取定义表示（$N$ 维），这个配分函数就是 Jones 多项式及其两变量推广，两个变量是 $N$ 与 $k$（解析延拓为复值）\src{355--356}；取 $G=SO(N)$、$R$ 为 $N$ 维表示则给出 Kauffman 多项式。

\textbf{这里要小心三件事}：
\begin{itemize}[leftmargin=1.6em]
\item 原文用的是 “we will claim” \src{355}，是断言，不是在这一节证明的；证明分散在 §2（弱耦合展开）与 §4（手术计算）。
\item \textbf{弱耦合极限是大 $k$}：$k\to\infty$ 时 $e^{\ii L}$ 剧烈振荡，鞍点法可用。这与 Yang--Mills 的 $1/g^2$ 展开方向相反，不能混为一谈。
\item 两变量里的 $k$ 在最终公式中要换成 $k+N$（或 $k+h^\vee$）——半形式修正，本文 §3 略去，见第 \ref{sec:canonical} 节。
\end{itemize}

\begin{boundary}
\begin{itemize}[leftmargin=1.4em]
\item 原文给出：(1.1)--(1.6)、$k$ 整性的必要性、$S^3$ 上复现 Jones 的断言。
\item 我们补的：变分公式 \eqref{eq:var} 的显式系数、$\Delta L=2\pi k m$ 的显式计算（原文只写 const$\cdot m$）。
\item 未证：$k$ 整性的“充分性”依赖测度与归一化的约定，原文未逐条论证；$S^3$ 上配分函数等于 Jones 多项式的严格证明不在本节。
\end{itemize}
\end{boundary}
